% CP-VI-0027
% target_chapter: VI/18 — Affine Schemes
% source_unit_id: IMP-DL-0B4CC237FA-P001
% source: Downloads/theory-of-algebraic-geometry-3.tex L177-L301
% solution_provenance: SOURCE_EXPLICIT_SOLUTION

\subsection*{Problem 11. The basic open subscheme \(D(f)\)}

Let \(A\) be a ring, let \(f\in A\), and let
\[
X=\operatorname{Spec} A.
\]
Show that, as locally ringed spaces,
\[
\bigl(D(f),\mathcal O_X|_{D(f)}\bigr)\cong \operatorname{Spec} A_f .
\]

\textbf{Solution.}

The statement says that passing from \(A\) to the localization \(A_f\) is geometrically the same as restricting \(\operatorname{Spec} A\) to the basic open set \(D(f)\).

Recall that
\[
D(f)=\{\mathfrak p\in \operatorname{Spec} A : f\notin \mathfrak p\}.
\]

We first compare the underlying topological spaces.

Prime ideals of \(A_f\) correspond exactly to prime ideals of \(A\) which do not contain \(f\). More precisely, we have mutually inverse maps
\[
\mathfrak p \longmapsto \mathfrak pA_f,
\]
and
\[
\mathfrak q \longmapsto \mathfrak q\cap A,
\]
where
\[
\mathfrak p\in \operatorname{Spec} A,\qquad f\notin \mathfrak p,
\]
and
\[
\mathfrak q\in \operatorname{Spec} A_f.
\]
Therefore there is a bijection
\[
\operatorname{Spec} A_f \longrightarrow D(f),
\qquad
\mathfrak q\longmapsto \mathfrak q\cap A.
\]

We now show that this bijection is a homeomorphism.

A basic open subset of \(\operatorname{Spec} A_f\) has the form
\[
D\left(\frac g1\right)
\]
for some \(g\in A\). Under the above correspondence, this subset maps to
\[
D(g)\cap D(f)\subseteq \operatorname{Spec} A.
\]
But
\[
D(g)\cap D(f)=D(fg).
\]
Therefore basic opens in \(\operatorname{Spec} A_f\) correspond exactly to basic opens inside \(D(f)\). Hence the bijection
\[
\operatorname{Spec} A_f \cong D(f)
\]
is a homeomorphism.

It remains to compare the structure sheaves.

Let
\[
D(g)\cap D(f)=D(fg)
\]
be a basic open subset of \(D(f)\). On the restricted structure sheaf, we have
\[
\mathcal O_X|_{D(f)}(D(fg))
=
\mathcal O_X(D(fg))
=
A_{fg}.
\]
On the other hand, the corresponding basic open subset of \(\operatorname{Spec} A_f\) is
\[
D\left(\frac g1\right),
\]
and therefore
\[
\mathcal O_{\operatorname{Spec} A_f}\left(D\left(\frac g1\right)\right)
=
(A_f)_{g/1}.
\]
But localization is associative, so there is a canonical isomorphism
\[
(A_f)_{g/1}\cong A_{fg}.
\]
Explicitly,
\[
\frac{a/f^m}{(g/1)^n}
\longmapsto
\frac{a}{f^mg^n}.
\]
Thus the two sheaves agree on a basis of open sets.

Moreover, these identifications are compatible with restriction maps. Indeed, if
\[
D(h)\cap D(f)\subseteq D(g)\cap D(f),
\]
then the restriction maps are exactly the usual localization maps
\[
A_{fg}\longrightarrow A_{fh}.
\]
Therefore the sheaves are canonically isomorphic.

Finally, the stalks agree. If \(\mathfrak p\in D(f)\), then the corresponding prime of \(A_f\) is \(\mathfrak pA_f\), and
\[
(A_f)_{\mathfrak pA_f}\cong A_{\mathfrak p}.
\]
Hence the induced morphism on stalks is an isomorphism of local rings.

Therefore
\[
\boxed{
	\bigl(D(f),\mathcal O_X|_{D(f)}\bigr)\cong \operatorname{Spec} A_f
}
\]
as locally ringed spaces.
